

Cela résulte immédiatement de 1.9 et de la démonstration de 1.5.

\subsection*{2. Complexe cotangent relatif}

\textbf{2.1.} Nous donnerons ici une définition du complexe cotangent qui, sans avoir toute la généralité possible (voir Exp. 0 4.4), sera suffisante pour les besoins du théorème de Riemann--Roch; bien entendu, la définition donnée dans exp. 0 coïncide dans le cas envisagé avec la définition donnée ici.

Soit @@M0@@ un morphisme lissifiable, c'est-à-dire admettant une factorisation :
@@M1@@
où @@M2@@ est un morphisme lisse, et @@M3@@ une immersion qu'on peut supposer fermée. Soit @@M4@@ l'idéal de @@M5@@ dans @@M6@@. La dérivation canonique :
@@M7@@
définit un homomorphisme :
@@M8@@
qui s'annule sur @@M9@@, et définit donc un homomorphisme, encore noté @@M10@@ :
@@M11@@
On vérifie aisément que cet homomorphisme est @@M12@@-linéaire. On définit alors un complexe de chaînes @@M13@@ en posant :
@@M14@@
et en prenant comme homomorphisme @@M15@@ l'homomorphisme (2.1.1).

